\section*{§8. Распределение Максвелла}
\addcontentsline{toc}{section}{§8. Распределение Максвелла}

\begin{basics}
\sub{Что такое плотность вероятности $w$}
$w(x)\,\dd x$ — вероятность того, что величина у одной молекулы попадёт в интервал $(x,x+\dd x)$
(= средняя доля молекул с таким $x$). Число молекул: $\dd N=Nw(x)\dd x$.
Нормировка $\int w\,\dd x=1$. Среднее: $\avg{f}=\int f(x)w(x)\dd x$; дисперсия $\sigma^2=\avg{x^2}-\avg x^2$;
относительная флуктуация $\delta=\sigma/\avg x$.
Независимые величины: $w(x,y)=w(x)w(y)$.

\sub{Проекция скорости — гаусс}
\[w(v_x)=\sqrt{\frac{m}{2\pi kT}}\exp\Big(-\frac{mv_x^2}{2kT}\Big),\qquad \avg{v_x}=0,\quad\avg{v_x^2}=\frac{kT}{m}.\]
$m$ — масса \emph{одной} молекулы ($m=\mu/N_A$, $\frac{kT}{m}=\frac{RT}{\mu}$), $T$ — температура газа.
Проекции $v_x,v_y,v_z$ независимы.

\sub{Вектор скорости}
\[w(\vec v)=w(v_x)w(v_y)w(v_z)=\Big(\frac{m}{2\pi kT}\Big)^{3/2}\exp\Big(-\frac{mv^2}{2kT}\Big)\]
— плотность на единицу объёма \emph{пространства скоростей} $\dd^3v=\dd v_x\dd v_y\dd v_z$. Зависит только от $|\vec v|$: все направления равновероятны.

\sub{Модуль скорости (распределение Максвелла)}
Слой $(v,v+\dd v)$ в пространстве скоростей имеет объём $4\pi v^2\dd v$:
\[w(v)=4\pi\Big(\frac{m}{2\pi kT}\Big)^{3/2}v^2\exp\Big(-\frac{mv^2}{2kT}\Big).\]
Характерные скорости:
\begin{align*}
\vp&=\sqrt{2kT/m} &&\text{наивероятнейшая (максимум $w$)},\\
\vav&=\sqrt{8kT/\pi m} &&\text{средняя},\\
v_{\text{кв}}&=\sqrt{\avg{v^2}}=\sqrt{3kT/m} &&\text{среднеквадратичная};\qquad \vp<\vav<v_{\text{кв}}.
\end{align*}
Средняя кинетическая энергия $\avg E=\frac m2\avg{v^2}=\frac32kT$.

\sub{Замена переменной}
Если $y=y(x)$ монотонна: $w(y)=w(x(y))\,\Big|\dfrac{\dd x}{\dd y}\Big|$ (вероятность в соответствующих интервалах одна и та же).
Для нескольких переменных — на якобиан. Максимум $w(y)$ \emph{не} совпадает с $y(x^*)$.

\sub{Поток молекул на стенку (эффузия)}
Число ударов о единицу площади за единицу времени:
$j=n\int_0^\infty v_xw(v_x)\dd v_x=\frac14n\vav$.
Через малое отверстие площади $S$ вылетает $\frac14nS\vav$ молекул в секунду.

\sub{Гауссовы интегралы}
\[I_k=\int_0^\infty x^{k}e^{-ax^2}\dd x=\frac{\Gamma\big(\frac{k+1}2\big)}{2a^{(k+1)/2}},\qquad I_{k+2}=\frac{k+1}{2a}I_k:\]
\[I_0=\frac12\sqrt{\frac\pi a},\ I_1=\frac1{2a},\ I_2=\frac{\sqrt\pi}{4a^{3/2}},\ I_3=\frac1{2a^2},\ I_4=\frac{3\sqrt\pi}{8a^{5/2}},\ I_5=\frac1{a^3}.\]
$\int_{-\infty}^\infty=2I_k$ для чётных $k$ и $0$ для нечётных.
Функция ошибок $\erf x=\frac2{\sqrt\pi}\int_0^xe^{-t^2}\dd t$, $\erf\infty=1$.
\end{basics}

%-------------------------------------------------------------
\begin{problem}{8.1}{ex}{Оценить отношение числа молекул азота, проекция скорости которых на произвольную ось лежит в интервале
$(v_1-\Delta v_1,\,v_1+\Delta v_1)$, к числу молекул с проекцией в интервале $(v_2-\Delta v_2,\,v_2+\Delta v_2)$, при $T=300$\,К.
$\mu=28$\,г/моль, $v_1=300$\,м/с, $\Delta v_1=0{,}31$\,м/с, $v_2=500$\,м/с, $\Delta v_2=0{,}51$\,м/с.}
\textbf{1.} Интервалы узкие ($\Delta v\ll v$), поэтому число молекул в интервале $\approx$ плотность $\times$ ширина:
$\Delta N_i=N\,w(v_i)\cdot2\Delta v_i$, где $w(v_x)=\sqrt{\frac m{2\pi kT}}e^{-mv_x^2/2kT}$.

\textbf{2.} В отношении $N$ и предэкспоненциальный множитель сокращаются:
\[\frac{\Delta N_1}{\Delta N_2}=\frac{\Delta v_1}{\Delta v_2}\exp\Big[\frac{m(v_2^2-v_1^2)}{2kT}\Big].\]

\textbf{3. Числа:} $m=\mu/N_A=4{,}65\cdot10^{-26}$\,кг;\;
$\frac{m(v_2^2-v_1^2)}{2kT}=\frac{4{,}65\cdot10^{-26}\cdot1{,}6\cdot10^5}{2\cdot1{,}38\cdot10^{-23}\cdot300}=0{,}90$;\;
$\frac{0{,}31}{0{,}51}e^{0{,}90}=0{,}61\cdot2{,}46$.
\answer{$\Delta N_1/\Delta N_2\approx1{,}5$.}
\end{problem}

%-------------------------------------------------------------
\begin{problem}{8.2}{exdz}{Используя распределение Максвелла, найти плотность вероятности распределения молекул по кинетическим энергиям $w(E)$.}
\textbf{1. Замена переменной.} $E=\frac{mv^2}{2}\Rightarrow v=\sqrt{\frac{2E}{m}}$, $\frac{\dd v}{\dd E}=\frac{1}{\sqrt{2mE}}$.
Вероятность попасть в соответствующие интервалы одна и та же: $w(E)\dd E=w(v)\dd v$, т.е.\ $w(E)=w(v)\frac{\dd v}{\dd E}$.

\textbf{2. Подставляем} $w(v)=4\pi\big(\frac m{2\pi kT}\big)^{3/2}v^2e^{-mv^2/2kT}$, $v^2=\frac{2E}m$:
\[w(E)=4\pi\Big(\frac m{2\pi kT}\Big)^{3/2}\frac{2E}{m}e^{-E/kT}\frac{1}{\sqrt{2mE}}=\frac{2}{\sqrt\pi}(kT)^{-3/2}\sqrt E\,e^{-E/kT}.\]

\textbf{3. Характерные значения.} Максимум: $\frac{\dd}{\dd E}\big(\sqrt Ee^{-E/kT}\big)=\big(\frac{1}{2\sqrt E}-\frac{\sqrt E}{kT}\big)e^{-E/kT}=0\Rightarrow E^*=\frac{kT}2$.
Среднее: $\avg E=\frac m2\avg{v^2}=\frac m2\cdot\frac{3kT}m=\frac32kT$.
\answer{$w(E)=\dfrac{2}{\sqrt\pi}(kT)^{-3/2}\sqrt E\,e^{-E/kT}$,\; $E^*=\dfrac{kT}2$,\; $\avg E=\dfrac32kT$.}
\note{$E^*\ne\frac{m\vp^2}{2}=kT$ — наиболее вероятные значения при замене переменной не пересчитываются.}
\end{problem}

%-------------------------------------------------------------
\begin{problem}{8.3}{exdz}{Для идеального газа из $N$ частиц найти относительные флуктуации энергии одной молекулы $\delta(E)$ и всего газа $\delta(E_N)$.}
\textbf{Найти:} $\delta=\sigma/\text{среднее}$, где $\sigma^2=\avg{E^2}-\avg E^2$.

\textbf{1. Средние по $w(E)$ из 8.2.} Замена $E=kTt$ сводит интегралы к $\int_0^\infty t^{k}e^{-t}\dd t=\Gamma(k+1)$:
\[\avg E=\frac{2kT}{\sqrt\pi}\Gamma\big(\tfrac52\big)=\frac{2kT}{\sqrt\pi}\cdot\frac{3\sqrt\pi}{4}=\frac32kT,\]
\[\avg{E^2}=\frac{2(kT)^2}{\sqrt\pi}\Gamma\big(\tfrac72\big)=\frac{2(kT)^2}{\sqrt\pi}\cdot\frac{15\sqrt\pi}{8}=\frac{15}4(kT)^2.\]

\textbf{2. Одна молекула:}
\[\sigma^2=\frac{15}4(kT)^2-\frac94(kT)^2=\frac32(kT)^2,\qquad \delta(E)=\frac{\sqrt{3/2}\,kT}{\frac32kT}=\sqrt{\frac23}.\]

\textbf{3. Весь газ.} $E_N=E_1+\dots+E_N$, молекулы независимы: средние и дисперсии складываются,
$\avg{E_N}=N\avg E$, $\sigma_N^2=N\sigma^2$, поэтому $\delta(E_N)=\frac{\sqrt N\sigma}{N\avg E}=\frac{\delta(E)}{\sqrt N}$.
\answer{$\delta(E)=\sqrt{2/3}\approx0{,}82$,\; $\delta(E_N)=\sqrt{2/(3N)}$.}
\end{problem}

%-------------------------------------------------------------
\begin{problem}{8.4}{ex}{Найти вероятность того, что кинетическая энергия молекулы идеального газа не превышает $E_0$. Температура $T$.}
\textbf{1.} $P(E_0)=\int_0^{E_0}w(E)\dd E$ с $w(E)$ из 8.2. Замена $E=kTt^2$ ($\dd E=2kTt\,\dd t$), $x_0=\sqrt{E_0/kT}$:
\[P=\frac{2}{\sqrt\pi}\int_0^{x_0}t\,e^{-t^2}\cdot2t\,\dd t=\frac4{\sqrt\pi}\int_0^{x_0}t^2e^{-t^2}\dd t.\]

\textbf{2. По частям} ($u=t$, $\dd v=te^{-t^2}\dd t$, $v=-\frac12e^{-t^2}$):
$\int_0^{x_0}t^2e^{-t^2}\dd t=-\frac{x_0}2e^{-x_0^2}+\frac12\int_0^{x_0}e^{-t^2}\dd t=-\frac{x_0}2e^{-x_0^2}+\frac{\sqrt\pi}4\erf x_0$.
\answer{$P(E_0)=\erf\sqrt{y}-2\sqrt{y/\pi}\;e^{-y}$, где $y=E_0/kT$.}
\note{энергию меньше $E^*=kT/2$ имеют около 20\% молекул, меньше $\avg E=\frac32kT$ — около 60\%.}
\end{problem}

%-------------------------------------------------------------
\begin{problem}{8.5}{ex}{Жидкость при температуре $T$ разлита тонким слоем; движение молекул плёнки двумерно, проекции скоростей независимы
и распределены по Гауссу. Найти $w(E)$, наиболее вероятные и средние значения энергии и модуля импульса. Масса молекулы $m$.}
\textbf{1. Распределение по модулю импульса.} $w(p_x,p_y)=\frac{1}{2\pi mkT}e^{-(p_x^2+p_y^2)/2mkT}$. В полярных координатах
элемент площади $p\,\dd p\,\dd\varphi$, интеграл по $\varphi$ даёт $2\pi$:
\[w(p)=\frac{p}{mkT}e^{-p^2/2mkT}.\]
Максимум: $p^*=\sqrt{mkT}$. Среднее: $\bar p=\frac1{mkT}\int_0^\infty p^2e^{-p^2/2mkT}\dd p=\sqrt{\frac{\pi mkT}2}$.

\textbf{2. По энергии.} $E=\frac{p^2}{2m}$, $\frac{\dd p}{\dd E}=\frac mp$:\; $w(E)=w(p)\frac mp=\frac1{kT}e^{-E/kT}$.
Это убывающая экспонента: $E^*=0$, $\avg E=kT$.
\answer{$w(E)=\dfrac1{kT}e^{-E/kT}$;\; $E^*=0$, $\avg E=kT$;\; $p^*=\sqrt{mkT}$, $\bar p=\sqrt{\pi mkT/2}$.}
\end{problem}

%-------------------------------------------------------------
\begin{problem}{8.6}{exdz}{Найти при температуре $T$ среднее значение модуля скорости относительного движения двух молекул идеального газа. Масса молекулы $m$.}
\textbf{1. Совместное распределение.} Молекулы независимы: $w(\vec v_1,\vec v_2)\propto e^{-m(v_1^2+v_2^2)/2kT}$.

\textbf{2. Новые переменные:} скорость центра масс $\vec V=\frac{\vec v_1+\vec v_2}2$ и относительная $\vec u=\vec v_1-\vec v_2$.
Тогда $v_1^2+v_2^2=2V^2+\frac{u^2}{2}$, и
\[w\propto e^{-2mV^2/2kT}\cdot e^{-(m/2)u^2/2kT}.\]
Распределение по $\vec u$ — максвелловское, но с массой $m/2$ вместо $m$.

\textbf{3. Средний модуль} — формула $\vav=\sqrt{8kT/\pi m}$ с заменой $m\to m/2$.
\answer{$\avg{|\vec v_1-\vec v_2|}=\sqrt{\dfrac{16kT}{\pi m}}=4\sqrt{\dfrac{kT}{\pi m}}=\sqrt2\,\vav$.}
\end{problem}

%-------------------------------------------------------------
\begin{problem}{8.7}{exdz}{Найти среднее число молекул, ударяющихся за время $\tau$ о плоскую стенку площадью $S$. Температура $T$, масса молекулы $m$, концентрация $n$.}
\textbf{1.} Ось $x$ — перпендикулярно стенке, к ней. За время $\tau$ о стенку ударятся молекулы с проекцией $v_x>0$,
находящиеся в цилиндре с основанием $S$ и высотой $v_x\tau$. Таких молекул с проекцией в $(v_x,v_x+\dd v_x)$:
$\dd N=n\,S\,v_x\tau\cdot w(v_x)\dd v_x$.

\textbf{2. Суммируем по всем $v_x>0$:}
\[N=nS\tau\sqrt{\frac{m}{2\pi kT}}\int_0^\infty v_xe^{-mv_x^2/2kT}\dd v_x=nS\tau\sqrt{\frac{m}{2\pi kT}}\cdot\frac{kT}{m}=nS\tau\sqrt{\frac{kT}{2\pi m}}.\]
\answer{$N=nS\tau\sqrt{\dfrac{kT}{2\pi m}}=\dfrac14nS\tau\,\vav$.}
\end{problem}

%-------------------------------------------------------------
\begin{problem}{8.8}{exdz}{Считая концентрацию $n$ практически постоянной, найти средние модуль скорости и энергию молекул, вылетающих из малого
отверстия площадью $S$ в плоской стенке сосуда. Температура газа $T$, масса молекулы $m$.}
\textbf{1. Кто вылетает.} За время $\tau$ из молекул со скоростью $v$ под углом $\theta$ к нормали вылетят те, что в наклонном цилиндре объёмом
$Sv\tau\cos\theta$. Значит, число вылетевших со скоростью $v$ пропорционально $v\cdot w(v)$: быстрые молекулы вылетают чаще.

\textbf{2. Средние по вылетевшим} — усредняем с весом $v\,w(v)$:
\[\avg v_{\text{выл}}=\frac{\int v\cdot v\,w\,\dd v}{\int v\,w\,\dd v}=\frac{\avg{v^2}}{\vav},\qquad
\avg E_{\text{выл}}=\frac m2\frac{\int v^2\cdot v\,w\,\dd v}{\int v\,w\,\dd v}=\frac m2\frac{\avg{v^3}}{\vav}.\]

\textbf{3. Подставляем} $\avg{v^2}=\frac{3kT}m$, $\vav=\sqrt{\frac{8kT}{\pi m}}$, $\avg{v^3}=\frac{4}{\sqrt\pi}\big(\frac{2kT}m\big)^{3/2}$
(гауссов интеграл $I_5$):
\[\avg v_{\text{выл}}=\frac{3kT/m}{\sqrt{8kT/\pi m}}=3\sqrt{\frac{\pi kT}{8m}},\qquad
\frac{\avg{v^3}}{\vav}=\frac{4kT}{m}\ \Rightarrow\ \avg E_{\text{выл}}=2kT.\]
\answer{$\avg v=3\sqrt{\dfrac{\pi kT}{8m}}=\dfrac{3\pi}8\vav$,\; $\avg E=2kT$ (больше, чем $\frac32kT$ в сосуде).}
\end{problem}

%-------------------------------------------------------------
\begin{problem}{8.9}{exdz}{Тонкостенный пустой сосуд объёмом $V$ погружён в идеальный газ, занимающий гораздо больший объём; температура $T$ и
концентрация $n_0$ газа постоянны. В стенке сосуда образовалось малое отверстие площадью $S$. Найти зависимость концентрации газа в сосуде
от времени. Средний модуль скорости молекул внутри и снаружи одинаков и равен $\vav$.}
\textbf{1. Потоки через отверстие} (8.7): внутрь за время $\dd t$ влетает $\frac14n_0\vav S\,\dd t$, наружу вылетает $\frac14n\vav S\,\dd t$.

\textbf{2. Баланс числа молекул в сосуде} ($N=nV$):
\[V\frac{\dd n}{\dd t}=\frac14S\vav(n_0-n),\qquad n(0)=0.\]

\textbf{3. Решаем} разделением переменных: $\frac{\dd n}{n_0-n}=\frac{\dd t}{\tau}$, $\tau=\frac{4V}{S\vav}$ $\Rightarrow$
$-\ln\frac{n_0-n}{n_0}=\frac t\tau$.
\answer{$n(t)=n_0\big(1-e^{-t/\tau}\big)$,\; $\tau=\dfrac{4V}{S\vav}$.}
\note{пример: He, 300\,К ($\vav\approx1260$\,м/с), $V=1$\,л, $S=0{,}01$\,мм$^2$ $\Rightarrow\tau\approx320$\,с.}
\end{problem}

%-------------------------------------------------------------
\begin{problem}{8.10}{exdz}{Невесомый поршень площадью $S$ удерживается в пустом цилиндре пружиной жёсткостью $\kappa$. В начале поршень у левой стенки,
пружина не деформирована. В левой стенке делают малое отверстие площадью $s_0$. Через какое время $\tau$ давление на поршень станет
равным половине атмосферного $p_0$? $T=\text{const}$, воздух — идеальный газ, сжатие пружины квазистатическое, $\vav$ одинакова внутри и снаружи.}
\textbf{1. Связь числа молекул с давлением.} Поршень сдвинулся на $x$, газ занимает объём $Sx$. Равновесие поршня: $\kappa x=pS\Rightarrow x=\frac{pS}{\kappa}$.
Уравнение состояния: $NkT=p\cdot Sx=\frac{p^2S^2}{\kappa}$, откуда $kT\,\dd N=\frac{2S^2}{\kappa}p\,\dd p$.

\textbf{2. Поток через отверстие} (как в 8.9), концентрации $n=\frac p{kT}$:
$\frac{\dd N}{\dd t}=\frac14s_0\vav\frac{p_0-p}{kT}$.

\textbf{3. Приравниваем} $\dd N$ из п.\,1 и п.\,2 и разделяем переменные:
\[\dd t=\frac{8S^2}{\kappa s_0\vav}\cdot\frac{p\,\dd p}{p_0-p}.\]

\textbf{4. Интегрируем} от $p=0$ до $p_0/2$; $\frac{p}{p_0-p}=-1+\frac{p_0}{p_0-p}$:
\[\int_0^{p_0/2}\frac{p\,\dd p}{p_0-p}=\Big[-p-p_0\ln(p_0-p)\Big]_0^{p_0/2}=p_0\Big(\ln2-\frac12\Big).\]
\answer{$\tau=\dfrac{8p_0S^2}{\kappa s_0\vav}\Big(\ln2-\dfrac12\Big)$.}
\end{problem}

%-------------------------------------------------------------
\begin{problem}{8.11}{ex}{Найти давление равновесного излучения на идеально отражающие стенки, считая его газом фотонов, движущихся в произвольных
направлениях с одинаковым по модулю импульсом $\hbar\omega/c$. Концентрация фотонов $n$.}
\textbf{1. Сколько фотонов ударяется.} Все летят со скоростью $c$, направления равновероятны: доля с углом к нормали в $(\theta,\theta+\dd\theta)$
равна $\frac12\sin\theta\,\dd\theta$. Как в 8.7, о площадку $S$ за время $\tau$ ударится $\dd N=nSc\tau\cos\theta\cdot\frac12\sin\theta\,\dd\theta$.

\textbf{2. Импульс.} При зеркальном отражении фотон передаёт стенке $2\frac{\hbar\omega}{c}\cos\theta$.

\textbf{3. Полный импульс} за время $\tau$ равен $pS\tau$:
\[pS\tau=\int_0^{\pi/2}2\frac{\hbar\omega}{c}\cos\theta\cdot nSc\tau\cos\theta\cdot\frac12\sin\theta\,\dd\theta
=nS\tau\hbar\omega\int_0^{\pi/2}\cos^2\theta\sin\theta\,\dd\theta=\frac{nS\tau\hbar\omega}{3}.\]
\answer{$p=\dfrac{n\hbar\omega}{3}=\dfrac u3$, где $u=n\hbar\omega$ — плотность энергии излучения.}
\note{для одноатомного газа аналогично $p=\frac23u$.}
\end{problem}

%-------------------------------------------------------------
\begin{problem}{8.12}{ex}{В сосуде при температуре $T$ находятся $N_1$ молекул массой $m_1$ и $N_2$ молекул массой $m_2$. Найти плотность
распределения $\tilde w(v)$ всех молекул по скоростям.}
\textbf{1.} Число молекул с модулем скорости в $(v,v+\dd v)$ складывается из молекул обоих сортов:
$\dd N=\dd N_1+\dd N_2=N_1w_1(v)\dd v+N_2w_2(v)\dd v$, где $w_i$ — распределение Максвелла с массой $m_i$.

\textbf{2.} С другой стороны, $\dd N=(N_1+N_2)\tilde w(v)\dd v$. Сравниваем.
\answer{$\tilde w(v)=\dfrac{N_1w_1(v)+N_2w_2(v)}{N_1+N_2}$, $w_i(v)=4\pi\Big(\dfrac{m_i}{2\pi kT}\Big)^{3/2}v^2e^{-m_iv^2/2kT}$.}
\note{если массы сильно различаются — у графика два максимума.}
\end{problem}

%-------------------------------------------------------------
\begin{problem}{8.19}{dz}{Найти скорость $\tilde v$ из условия, что среднее число молекул со скоростями меньше $\tilde v$ равно половине их общего числа.
Масса молекулы $m$, температура $T$.}
\textbf{Найти:} $\tilde v$ из условия $\int_0^{\tilde v}w(v)\dd v=\frac12$.

\textbf{1. Безразмерная переменная} $x=v/\vp$, $\vp=\sqrt{2kT/m}$: тогда $\frac{mv^2}{2kT}=x^2$ и
\[\int_0^{\tilde v}w(v)\dd v=\frac4{\sqrt\pi}\int_0^{x_0}x^2e^{-x^2}\dd x,\qquad x_0=\tilde v/\vp.\]

\textbf{2. По частям} (как в 8.4): $\int_0^{x_0}x^2e^{-x^2}\dd x=-\frac{x_0}2e^{-x_0^2}+\frac{\sqrt\pi}4\erf x_0$, поэтому условие:
\[\erf x_0-\frac2{\sqrt\pi}x_0e^{-x_0^2}=\frac12.\]

\textbf{3.} Уравнение трансцендентное, решаем численно (подбором по таблице $\erf$): $x_0\approx1{,}088$.
\answer{$\tilde v=x_0\sqrt{2kT/m}\approx1{,}088\sqrt{2kT/m}$.}
\note{$\vp<\tilde v<\vav<v_{\text{кв}}$: в единицах $\vp$ это $1<1{,}088<1{,}128<1{,}225$.}
\end{problem}

%-------------------------------------------------------------
\begin{problem}{8.20}{dz}{Вычислить средние значения проекции и модуля проекции скорости молекулы идеального газа на некоторое направление.
Масса молекулы $m$, температура $T$.}
\textbf{1. Среднее $v_x$.} $w(v_x)$ — чётная функция, а $v_x$ — нечётная, интеграл по симметричному промежутку равен нулю: $\avg{v_x}=0$.

\textbf{2. Среднее $|v_x|$.} Подынтегральная функция чётная — берём удвоенный интеграл по $v_x>0$:
\[\avg{|v_x|}=2\sqrt{\frac m{2\pi kT}}\int_0^\infty v_xe^{-mv_x^2/2kT}\dd v_x=2\sqrt{\frac m{2\pi kT}}\cdot\frac{kT}{m}=\sqrt{\frac{2kT}{\pi m}}\]
(использовали $\int_0^\infty xe^{-ax^2}\dd x=\frac1{2a}$, $a=\frac m{2kT}$).
\answer{$\avg{v_x}=0$,\; $\avg{|v_x|}=\sqrt{\dfrac{2kT}{\pi m}}=\dfrac{\vav}2$.}
\end{problem}

%-------------------------------------------------------------
\begin{problem}{8.21}{dz}{В сосуде при температуре $T$ находится $N$ молекул идеального газа. Найти число частиц, у которых одновременно проекция
скорости на ось $Oz$ лежит в $(v_z,v_z+\dd v_z)$, а модуль перпендикулярной составляющей скорости — в $(v_\perp,v_\perp+\dd v_\perp)$. Масса молекулы $m$.}
\textbf{1. Исходное распределение} по вектору скорости:
$\dd N=N\big(\frac m{2\pi kT}\big)^{3/2}e^{-m(v_x^2+v_y^2+v_z^2)/2kT}\dd v_x\dd v_y\dd v_z$.

\textbf{2. Цилиндрические координаты} в пространстве скоростей: $v_x=v_\perp\cos\varphi$, $v_y=v_\perp\sin\varphi$;
$v_x^2+v_y^2=v_\perp^2$, $\dd v_x\dd v_y=v_\perp\dd v_\perp\dd\varphi$.

\textbf{3. Направление $\vec v_\perp$ не важно} — интегрируем по $\varphi$ от 0 до $2\pi$, получаем множитель $2\pi$.
\answer{$\dd N=2\pi N\Big(\dfrac m{2\pi kT}\Big)^{3/2}v_\perp\exp\Big[-\dfrac{m(v_z^2+v_\perp^2)}{2kT}\Big]\dd v_z\,\dd v_\perp$.}
\end{problem}

%-------------------------------------------------------------
\begin{problem}{8.23}{dz}{Проекции скорости молекул идеального газа на оси статистически независимы и распределены по Гауссу. Найти совместную
$w(v_x,v)$ и условную $w(v_x\,|\,v)$ плотности вероятности. Масса молекулы $m$, температура $T$.}
\textbf{1. Условная плотность.} Если модуль $v$ задан, то направление скорости равновероятно по всем направлениям.
Для равномерного распределения на сфере $\cos\theta$ равномерно распределён на $[-1,1]$
(площадь пояса сферы $2\pi\sin\theta\,\dd\theta=2\pi\,|\dd\cos\theta|$). Так как $v_x=v\cos\theta$, то $v_x$ равномерно распределена на $[-v,v]$:
\[w(v_x\,|\,v)=\frac{1}{2v},\qquad|v_x|\le v.\]

\textbf{2. Совместная плотность} = плотность $v$ $\times$ условная: $w(v_x,v)=w(v)\cdot w(v_x\,|\,v)$, где $w(v)$ — распределение Максвелла:
\[w(v_x,v)=4\pi\Big(\frac m{2\pi kT}\Big)^{3/2}v^2e^{-mv^2/2kT}\cdot\frac1{2v}.\]
\answer{$w(v_x\,|\,v)=\dfrac1{2v}$,\; $w(v_x,v)=2\pi\Big(\dfrac m{2\pi kT}\Big)^{3/2}v\,e^{-mv^2/2kT}$\; (при $|v_x|\le v$, иначе 0).}
\end{problem}
